\section{Dynamics and the Radial Wave Equation}

We formally solve the Time-Dependent Schrödinger Equation (TDSE) by projecting
the dynamics onto the eigenbasis of the Complete Set of Commuting Observables
(CSCO) $\{ \hat{H}, \hat{L}^2, \hat{L}_z \}$. This procedure separates the
trivial temporal evolution from the geometric spatial structure.

\subsection{Spectral Evolution}

The state vector $\ket{\psi(t)}$ is expanded in the continuous basis of
simultaneous eigenstates $\ket{k, \ell, m}$, normalized to the radial measure
$k^2 \mathrm{d}k$:
\begin{equation}
	\ket{\psi(t)} = \sum_{\ell=0}^{\infty} \sum_{m=-\ell}^{\ell}
	\int_{0}^{\infty} \mathrm{d}k \, k^2 \, c_{k\ell m}(t) \, \ket{k, \ell, m}.
\end{equation}
Substituting this expansion into the TDSE $i\hbar \partial_t \ket{\psi} =
\hat{H}\ket{\psi}$ and exploiting the orthogonality of the basis, the partial
differential equation decouples into an infinite set of independent ordinary
differential equations for the coefficients $c_{k\ell m}(t)$:
\begin{equation}
	i\hbar \dot{c}_{k\ell m}(t) = E_k c_{k\ell m}(t), \quad \text{where } E_k =
	\frac{\hbar^2 k^2}{2m}.
\end{equation}
The solution is exact and demonstrates that the angular quantum numbers are
constants of motion:
\begin{equation}
	c_{k\ell m}(t) = c_{k\ell m}(0) \exp\left( -\frac{i}{\hbar} \frac{\hbar^2
	k^2}{2m} t \right).
\end{equation}

\subsection{The Coordinate Representation (Radial Equation)}

To determine the spatial probability distribution, we must find the coordinate
representation of the basis kets, $\braket{\mathbf{r} | k, \ell, m}$. We exploit
the separation of variables in spherical coordinates:
\begin{equation}
	\braket{\mathbf{r} | k, \ell, m} = R_{k\ell}(r) Y_{\ell m}(\theta, \phi).
\end{equation}
Here, $Y_{\ell m}$ are the Spherical Harmonics derived previously. The radial
function $R_{k\ell}(r)$ is determined by the energy eigenvalue equation in
position space:
\begin{equation}
	\braket{\mathbf{r} | \hat{H} | k, \ell, m} = E_k \braket{\mathbf{r} |k,
	\ell, m}.
\end{equation}
Using the Laplacian decomposition $\nabla^2 = \partial_r^2 + \frac{2}{r}
\partial_r - \frac{\hat{L}^2}{\hbar^2 r^2}$, the equation becomes:
\begin{equation}
	-\frac{\hbar^2}{2m} \left[ \frac{1}{r^2} \frac{\partial}{\partial r} \left(
		r^2 \frac{\partial R_{k\ell}}{\partial r} \right) -
	\frac{\ell(\ell+1)}{r^2} R_{k\ell}(r) \right] = \frac{\hbar^2 k^2}{2m}
	R_{k\ell}(r).
\end{equation}

\subsubsection{The Spherical Bessel Equation}
Rearranging terms and introducing the dimensionless variable $\rho = kr$, we
obtain the Spherical Bessel differential equation:
\begin{equation}
	\frac{d^2 R}{d\rho^2} + \frac{2}{\rho} \frac{dR}{d\rho} + \left[ 1 -
	\frac{\ell(\ell+1)}{\rho^2} \right] R(\rho) = 0.
\end{equation}
The term $\frac{\ell(\ell+1)}{\rho^2}$ represents the centrifugal potential
barrier arising from angular momentum conservation, which prevents high-$\ell$
particles from penetrating near the origin.

\subsubsection{Free Spherical Waves}
The physically admissible solutions (regular at the origin $r=0$) are the
Spherical Bessel Functions of the first kind, denoted $j_\ell(kr)$.

\begin{equation}
	R_{k\ell}(r) = \sqrt{\frac{2}{\pi}} j_\ell(kr).
\end{equation}
(The normalization factor $\sqrt{2/\pi}$ is chosen to satisfy the delta-function
normalization $\int r^2 dr R_{k\ell} R_{k'\ell} = \delta(k-k')/k^2$).

\subsection{The Full Wavefunction Solution}

Combining the spectral evolution with the coordinate basis, the general solution
for the free particle in the angular momentum representation is:
\begin{equation}
	\psi(\mathbf{r}, t) = \sum_{\ell=0}^{\infty} \sum_{m=-\ell}^{\ell}
	\int_{0}^{\infty} \mathrm{d}k \, k^2 \, \underbrace{c_{k\ell m}(0) e^{-i
		E_k t/\hbar}}_{\text{Dynamics}} \, \underbrace{\sqrt{\frac{2}{\pi}}
		j_\ell(kr) Y_{\ell m}(\theta, \phi)}_{\text{Spatial Mode}}.
\end{equation}
This expansion represents the free particle as a superposition of spherical
standing waves, providing the natural basis for scattering problems where the
asymptotic behavior is defined by incoming and outgoing spherical states.

\section{Time Evolution and General Solution}

The diagonalization of the Hamiltonian within the CSCO basis $\{ \ket{k, \ell, m}
\}$ decouples the global evolution of the field $\psi(\mathbf{r}, t)$ into an
infinite set of independent dynamical modes. We construct the general solution
to the initial value problem by applying the unitary time-evolution operator to
the spectral expansion.

\subsection{Unitary Evolution of Spectral Modes}

The equation of motion derived implies
that each spectral amplitude $\psi_{k\ell m}(t)$ undergoes simple harmonic
oscillation in the complex plane. The formal solution is obtained via the action
of the propagator $\hat{U}(t) = \exp(-i\hat{H}t/\hbar)$:
\begin{equation}
	\psi_{k\ell m}(t) = \bra{k, \ell, m} \hat{U}(t) \ket{\psi(0)} = e^{-i E_k t
	/ \hbar} \psi_{k\ell m}(0).
\end{equation}

Explicitly, with the free particle dispersion relation $E_k = \frac{\hbar^2
k^2}{2m}$:
\begin{equation} \label{eq:mode_solution}
	\psi_{k\ell m}(t) = \psi_{k\ell m}(0) \exp\left( -i \frac{\hbar k^2}{2m} t
	\right).
\end{equation}
This result highlights the stationarity of the angular momentum eigenstates; the
probability density $|\psi_{k\ell m}(t)|^2$ is time-invariant, ensuring that
information encoded in the angular channels is conserved.

\subsection{Synthesis of the Dynamical State}

The full dynamical state $\ket{\psi(t)}$ is reconstructed by synthesizing the
evolved modes. Substituting the spectral solution (\ref{eq:mode_solution}) into
the basis expansion yields the general solution in abstract notation:
\begin{equation} \label{eq:general_solution_abstract}
	\ket{\psi(t)} = \sum_{\ell=0}^{\infty} \sum_{m=-\ell}^{\ell}
	\int_{0}^{\infty} \mathrm{d}k \, k^2 \, \psi_{k\ell m}(0) \, e^{-i
	\frac{\hbar k^2}{2m} t} \, \ket{k, \ell, m}.
\end{equation}

Projecting this state onto the coordinate basis $\bra{\mathbf{r}}$ and utilizing
the free spherical wave functions $\braket{\mathbf{r}}{k,\ell,m} =
\sqrt{\frac{2}{\pi}} j_\ell(kr) Y_{\ell m}(\Omega)$, we obtain the explicit
position-space wavefunction:
\begin{equation} \label{eq:general_solution_spatial}
	\psi(\mathbf{r}, t) = \sqrt{\frac{2}{\pi}} \sum_{\ell, m} Y_{\ell m}(\Omega)
	\int_{0}^{\infty} \mathrm{d}k \, k^2 \, \psi_{k\ell m}(0) \, j_\ell(kr) \,
	e^{-i \frac{\hbar k^2}{2m} t}.
\end{equation}
This integral transform constitutes the complete solution to the free-particle
Schrödinger equation in spherical coordinates.

\subsection{Physical Interpretation}

The structure of Eq. (\ref{eq:general_solution_spatial}) elucidates the
physical mechanism of free evolution:
\begin{itemize}
	\item \textbf{Angular Stability:} The quantum numbers $\ell$ and $m$ appear
		strictly as summation indices for the stationary harmonics $Y_{\ell m}$. The
		absence of time dependence in the angular sector confirms that the
		probability distribution over the orbital angular momentum channels is
		invariant, a direct consequence of the isotropy of free space ($[\hat{H},
		\hat{\mathbf{L}}] = 0$).
	\item \textbf{Radial Dispersion:} Time dependence is confined to the phase
		modulation of the radial modes. Since the dispersion relation is non-linear
		($E_k \propto k^2$), the group velocity $v_g = \partial_k \omega = \hbar
		k/m$ varies with wavenumber. This gradient in propagation speed across
		different $k$-components drives the characteristic wave packet spreading
		(dispersion) observed in the position representation.
\end{itemize}
Ultimately, the angular momentum representation resolves the complex diffusive
dynamics of the free particle into a static rotation of independent spectral
components in the complex plane.
